\subsection{SRIF Topics}
\begin{frame}
  \onslide<1->{
	\frametitle{Square Root Information Filters} 
	\framesubtitle{Topics}
}
\begin{itemize}
  \onslide<2->{
\item  Information Filters 
}
  \onslide<3->{
\item  Square Root Filters 
}
  \onslide<4->{
\item  Square Root Information Filters (SRIF) 
}
\end{itemize}
\end{frame}


\subsection{Information Filters}
\begin{frame}
  \onslide<1->{
	\frametitle{Square Root Information Filters} 
	\framesubtitle{Information Filters}
}
\begin{itemize}
  \onslide<2->{
\item Standard KF is sometimes called the Covariance Kalman Filter because it is formulated
using the covariance matrix $P$.  The Information Filter is more efficient when there are very large batches of 
measurements to process.
}

  \onslide<3->{
\item Information Matrix $\Lambda$:  Inverse of the Estimation Error Covariance Matrix $P$  \\
}
  \onslide<4->{
\item Recall $\gencovmatdef$  
}
  \onslide<5->{
\item  We define $\Lambda:=P^{-1}$ 
}
  \onslide<6->{
\item When $\Lambda$ is "large" $P$ is "small" (in norm)  
}
\end{itemize}
\end{frame}
\begin{frame}
  \onslide<1->{
	\frametitle{Square Root Information Filters} 
	\framesubtitle{Information Filters}
}
\begin{itemize}
  \onslide<2->{
\item  Consider the covariance update equation $\gencovupdateeq$ 
}
  \onslide<3->{
  % Taking inverses yields:  $\left[P^{+}\right]^{-1}=\left[P^{-}\right]^{-1}\left[ I-KH\right]^{-1}$ 
\item  Substitute the Kalman gain formula $\genkalgaineq$ into the above equation to yield\\
$P^{+}=P^{-}-P^{-}H\trans\left[HP^{-}H\trans+R\right]^{-1}HP^{-}$  
}
  \onslide<4->{
\item   To futher expand and then simplify, we need to apply the Matrix Inversion Lemma (also known as the Woodbury formula):\\
$\woodbury{A}{B}{C}{D}$
}
  \onslide<5->{
\item We obtain:  $\left[P^{+}\right]^{-1}=\left[P^{-}\right]^{-1}+H\trans R^{-1}H$, or,\\ in terms of $\Lambda$:
$\Lambda^{+}=\Lambda^{-}+H\trans R^{-1}H$  
}
\end{itemize}
\end{frame}

\begin{frame}
  \onslide<1->{
	\frametitle{Square Root Information Filters} 
	\framesubtitle{Information Filters}
}
\begin{itemize}
  \onslide<2->{
\item  Note that the information matrix is updated without knowledge of the gain!
}
  \onslide<3->{
\item  The gain comutation takes on a very simple form:  $K=\Lambda^{+}H\trans R^{-1}$  
}
  \onslide<4->{
\item  The state is then updated using the gain in the normal way 
}
  \onslide<5->{
\item  Standard KF requires inversion of $r\times r$ matrix ($r=$number of measurements); information filter
requires inversion of a couple of $n\times n$ inversions.   \\
Thus, if $r>>n$, it is more efficient to use the information filter.
}
\end{itemize}
\end{frame}
\begin{frame}
  \onslide<1->{
	\frametitle{Square Root Information Filters} 
	\framesubtitle{Information Filters}
}
\begin{itemize}
  \onslide<2->{
\item  The information matrix extrapolation equation is derived in a similar fashion by 
transforming the standard covariance extrapolation equation using once again the Woodbury formula
}
  \onslide<3->{
\item  This yields:\\ $\Lambda^{-}(k)=Q^{-1}-Q^{-1}\Phi\left[ \Lambda^{+}(k-1)+\Phi\trans Q^{-1} \Phi\right]^{-1}\Phi\trans Q^{-1}$ 
}
  \onslide<4->{
\item   If the process noise covariance $Q$ and/or the measurement noise covariance are constant (
or change only infrequently) with time, then the indicated inverses can be computed once at startup or perhaps only 
occasionaly.
}
\end{itemize}
\end{frame}

\subsection{Square Root Filters}
\begin{frame}
  \onslide<1->{
	\frametitle{Square Root Information Filters} 
	\framesubtitle{Square Root Filters}
}
\begin{itemize}
  \onslide<2->{
\item Motivation: Covariance matrices are supposed to be nonsingular, positive semi-definite; they are invertible  
}
  \onslide<3->{
\item  But due to differences in uncertainy values for different states, they may be ill-conditioned for inversion 
}
  \onslide<4->{
\item  Conditioning refers to the condition number of a matrix:  $\kappa (P):=\sigma_{\text{max}}/\sigma_{\text{min}}$,\\
where $\sigma_{\text{max}}$ is the maximum singular value of $P$ and $\sigma_{\text{min}}$ is the minimum singular value.\\
\begin{footnotesize} 
Singular values of real-valued $P$ are square roots of eigenvalues of $P\trans P$ (or $PP\trans$)
\end{footnotesize}
}
  \onslide<5->{
\item  Note that $\kappa(P)\geq 1$, for a nonsingular matrix $P$, since all of its singular values are positive
}
\end{itemize}
\end{frame}

\begin{frame}
  \onslide<1->{
	\frametitle{Square Root Information Filters} 
	\framesubtitle{Square Root Filters}
}
\begin{itemize}
  \onslide<2->{
\item Example of 2-state model: Highly uncertain first state, but highly certain second state:\\
$P=
\begin{bmatrix}
10^{10} & 0\\
0 & 10^{-10}
\end{bmatrix}$  
}
  \onslide<3->{
\item $\kappa(P)=10^{10}/10^{-10}=10^{20}$, since singular values of a diagonal matrix 
are absolute values of the diagonal elements  
}
  \onslide<4->{
\item Although $P$ is technically nonsingular, on an embedded digital computer with limited word length 
it may become singular   
}
  \onslide<5->{
\item  MATLAB Warning: Matrix is close to singular or badly scaled. Results may be inaccurate. RCOND =
1.000000e-20 
}
\end{itemize}
\end{frame}

\begin{frame}
  \onslide<1->{
	\frametitle{Square Root Information Filters} 
	\framesubtitle{Square Root Filters}
}
\begin{itemize}
  \onslide<2->{
\item If we can factor $P$ as the product of a matrix and its transpose, i.e., $P=SS\trans$, then the matrix $S$  
is said to be the matrix square root of $P$\\
\begin{footnotesize}Note: There are several alternative definitions in the literature $SS\trans,S^2,S\trans S$\end{footnotesize}
}
  \onslide<3->{
\item Key idea:  if $P=SS\trans$, then $\kappa(S)=\sqrt{\kappa(P)}$  
}
  \onslide<3->{
\item Matrix square roots are not unique  
}
  \onslide<4->{
\item But all square roots of a matrix are related by orthogonal transformations (i.e., 
rotations in n-dimensional space)\\
$T$ is \textbf{orthogonal} if its transpose is its inverse:   $TT\trans = T\trans T=I$  
}
  \onslide<5->{
\item If $P=S_1 S_1\trans=S_2 S_2\trans$, then there is an orthogonal matrix $T$ such that $S_2=S_1 T$
}
\end{itemize}
\end{frame}

\begin{frame}
  \onslide<1->{
	\frametitle{Square Root Information Filters} 
	\framesubtitle{Square Root Filters: Orthogonal Matrices}
}
\begin{itemize}
  \onslide<2->{
\item Important Properties of Orthogonal Matrices:  
}
\begin{itemize}
  \onslide<3->{
\item Length Invariance (Isometry):  $\norm{Tx}=\norm{x}$  \\
where $\norm{x}=\sqrt{x\trans x}$ is Euclidean norm\\ (also preserves dot product between two vectors)
}
  \onslide<4->{
\item Angle Invariance:  The angle between the vectors $x$ and $y$ is the same as that between $Tx$ and $Ty$ 
}
  \onslide<5->{
\item  Standard Gaussian Invariance: If $v$ is zero-mean Gaussian random vector such that $\expect{vv\trans}=I$, 
then for orthogonal $T$ the 
vector $Tv$ is also zero-mean Gaussian and $\expect{Tv(Tv)\trans}=I$
}
\end{itemize}
  \onslide<6->{
    \item Length and Standard Gaussian Invariances play crucial role in square root filtering
}
\end{itemize}
\end{frame}

\subsection{Bierman Derivation of SRIF}
\begin{frame}
  \onslide<1->{
	\frametitle{Square Root Information Filters} 
	\framesubtitle{Least Squares Derivation of SRIF (Bierman)}
}
\begin{itemize}
  \onslide<2->{
\item Consider measurement model:  $z=Hx+v$  
}
  \onslide<3->{
\item  Least squares solution minimizes means square error:  $$J=\sum_{j=1}^m v^2(j)=v\trans v$$ 
}
  \onslide<4->{
\item As a function of $x$ this becomes $$J(x)=(z-Hx)\trans (z-Hx)$$  
}
\end{itemize}
\end{frame}

\begin{frame}
  \onslide<1->{
	\frametitle{Square Root Information Filters} 
	\framesubtitle{Least Squares Derivation of SRIF (Bierman)}
}
\begin{itemize}
  \onslide<2->{
\item  Differentiating $J(x)$ with respect to $x$ (taking into account non-commutative multiplication of matrices),
we obtain the least squares \textbf{normal equations}: $$H\trans H x=H\trans z$$
}
  \onslide<3->{
\item  The least squares solution $x_{\text{ls}}$ is therefore $$x_{\text{ls}}=(H\trans H)^{-1}H\trans z$$ 
}
\end{itemize}
\end{frame}

\begin{frame}
  \onslide<1->{
	\frametitle{Square Root Information Filters} 
	\framesubtitle{Least Squares Derivation of SRIF (Bierman)}
}
\begin{itemize}
  \onslide<2->{
\item   Interpreting the errors $v$ statistically and assuming both $\expect{v}=0$ and $\expect{vv\trans}=I$,
we derive $$H\trans H \left(x_{\text{ls}}-x\right)=H\trans v$$
}
  \onslide<3->{
\item  And it can be shown that $$P_{\text{ls}}=(H\trans H)^{-1}$$ 
}
  \onslide<4->{
\item  Taking the matrix inverse of both sides reveals SRI matrix $$\Lambda_{\text{ls}} = H\trans H$$ 
}
\end{itemize}
\end{frame}

\begin{frame}
  \onslide<1->{
	\frametitle{Square Root Information Filters} 
	\framesubtitle{Least Squares Derivation of SRIF (Bierman)}
}
\begin{itemize}
  \onslide<2->{
\item  The above derivation was \textit{a posteriori}, i.e., after the measurement
}
  \onslide<3->{
\item  We can now incorpoate the \textit{a priori} estimate $\hat{x}^{-}$ and SRI matrix $\Lambda^{-}$
}
  \onslide<4->{
\item  The previous performance functional is now modified to become
$$
J_1(x):=(x-\hat{x}^{-})\trans \Lambda^{-} (x-\hat{x}^{-})+(z-Hx)\trans (z-Hx)
$$
}
\end{itemize}
\end{frame}

\begin{frame}
  \onslide<1->{
	\frametitle{Square Root Information Filters} 
	\framesubtitle{Least Squares Derivation of SRIF (Bierman)}
}
\begin{itemize}
  \onslide<2->{
\item  The vector that minimizes $J_1$ is denoted $\hat{x}^{+}_{\text{ls}}$ and is easily shown to satisfy
$$
(\Lambda^{-}+H\trans H)\hat{x}^{+}_{\text{ls}}=\Lambda^{-}\hat{x}^{-}+H\trans z
$$

}
  \onslide<3->{
\item   Finally, the updated covariance is
$$
P^{+}_{\text{ls}}=(\Lambda^{-}+H\trans H)^{-1}
$$
}
  \onslide<4->{
\item  Bierman points out that $\Lambda^{-}$ need not be non-singular
}
  \onslide<5->{
\item  In fact, one of the advantages of using the information matrix instead of the covariance 
is that many entries of $\Lambda^{-}$ may be 0 (meaning no information=infinite uncerrtainty); we 
may even start with $\Lambda^{-}=0$
}
\end{itemize}
\end{frame}

\begin{frame}
  \onslide<1->{
	\frametitle{Square Root Information Filters} 
	\framesubtitle{SRIF Recursion (Bierman)}
}
\begin{itemize}
  \onslide<2->{
\item Factor the \textit{a priori} $\Lambda^{-}$ into square root factors:  
$\Lambda^{-}=\left[S^{-}\right]\trans S^{-}$
}
  \onslide<3->{
\item This allows us to rewrite first term of $J_1(x)$ as
$$
(x-\hat{x}^{-})\trans \Lambda^{-} (x-\hat{x}^{-})=(z^{-}-S^{-}x)\trans (z^{-}-S^{-}x)
$$
where $z^{-}= S^{-}x^{-}$
}
  \onslide<4->{
\item  We can now treat the \textit{a priori} estimate and covariance information 
as if they were an additional (pseudo) measurement and augment the least squares problem
$$
\begin{bmatrix}
z^{-}\\
z
\end{bmatrix}
=
\begin{bmatrix}
S^{-}\\
H
\end{bmatrix}
x +
\begin{bmatrix}
v^{-}\\
v
\end{bmatrix}
$$
}
\end{itemize}
\end{frame}

\begin{frame}
  \onslide<1->{
	\frametitle{Square Root Information Filters} 
	\framesubtitle{Summary}
}
\begin{itemize}
  \onslide<2->{
\item  Bierman:  "The interpretation of the \textit{a priori} as additional observations is the cornerstone
of the square root information filter." 
}
  \onslide<3->{
\item  Results do not depend on which square root is used;  facilitates use of orthogonal transformations 
(Householder, Givens matrices) to put the square roots into numerical convenient forms 
}
  \onslide<4->{
\item   Bierman and Thornton (JPL) combined this SRIF formulation with Cholesky decomposition to devise the U-D 
factorization
}
\end{itemize}
\end{frame}
\begin{frame}
  \onslide<1->{
	\frametitle{UD Factorization} 
	\framesubtitle{Bierman and Thornton}
}
\begin{itemize}
  \onslide<2->{
  \item Based on Cholesky Decomposition
}
  \onslide<3->{
  \item Unit Triangular Matrices: $U$ has 1's on the diagonal and all 0's below the diagonal.  E.g.,
\[
U=\begin{bmatrix}
1 & 5 & 7 \\
0 & 1 & 3\\
0 & 0 & 1
\end{bmatrix}
\]
}
\end{itemize}
\end{frame}

\begin{frame}
  \onslide<1->{
	\frametitle{UD Factorization} 
	\framesubtitle{}
}
\begin{itemize}
  \onslide<2->{
  \item Factor $P^{-}$ as $UDU\trans$, where $U$ is unit upper triangular and $D$ is diagonal.
}
  \onslide<3->{
  \item This then allows us to "invert" $P^{-}$ by solving the equation unit upper triangular equation $UDU\trans X=H$ for the matrix $X$, which will yield $X=\left[HPH\trans + R\right]^{-1}H$
}
  \onslide<4->{
  \item This last expression is used on the covariance update equation in $KH$
}
\end{itemize}
\end{frame}


